\section{Верификация и надграждаемост чрез нови конектори}
\label{sec:verification_extensibility}

Една от централните претенции на разработката е, че библиотеката е едновременно верифицируема и надграждаема. Това изисква не просто наличие на тестове, а ясно дефиниран connector contract, който позволява нов backend да бъде добавен без промяна в query API и без разпадане на compile-time гаранциите.

В предходните глави беше показано, че библиотеката вече работи с релационни и нерелационни backend-и. Настоящата глава формализира как това е възможно и как трябва да се добавя следващ backend. С други думи, тук надграждаемостта престава да бъде общо качество на дизайна и се превръща в методология с конкретни изисквания, signatures, тестови критерии и нива на завършеност.

\subsection{Формален contract на конекторния слой}

Конекторът се описва чрез тип \code{DB} и специализация \code{connector\_trait<DB>}. Това е структурен, а не наследствен contract. Няма общ виртуален интерфейс \code{IConnector}. Причината е принципна: различните backend-и имат различни възможности и ограничения, които трябва да бъдат видими на compile-time. Виртуалната йерархия би изравнила твърде много различия и би прехвърлила част от проверките към runtime.

Минималният работещ конектор трябва да осигури \code{wire\_type}, \code{cursor\_type} и подходящи \code{execute(...)} entry points. Допълнителните способности се означават чрез capability tags. Така contract-ът остава едновременно строг и разширяем. В тази конструкция няма ``магически'' registration слой: компилаторът вижда специализацията директно и може да провери дали тя удовлетворява \code{is\_connector} concept-а.

\begin{figure}[H]
\centering
\begin{tikzpicture}[node distance=1.35cm and 1.2cm]
    \node[ormaccent] (dbtype) {\textbf{Нов backend тип}\\resource wrapper + lifecycle};
    \node[ormblock, below=of dbtype] (trait) {\code{connector\_trait<DB>}\\mandatory signatures};
    \node[ormblock, below left=of trait] (caps) {Capability\\декларации};
    \node[ormblock, below right=of trait] (exec) {\code{execute(...)}\\materialization path};
    \node[ormblock, below=of $(caps)!0.5!(exec)$] (tests) {Unit + integration\\верификация};

    \draw[ormline] (dbtype) -- (trait);
    \draw[ormline] (trait) -- (caps);
    \draw[ormline] (trait) -- (exec);
    \draw[ormline] (caps) -- (tests);
    \draw[ormline] (exec) -- (tests);
\end{tikzpicture}
\caption{Архитектурен workflow за добавяне на нов connector}
\label{fig:connector_workflow}
\end{figure}

Фигура~\ref{fig:connector_workflow} показва, че добавянето на backend не е едноетапно действие. То започва с ясно дефиниран resource wrapper, продължава с trait contract, capability профил и execution materialization, и завършва едва след тестова верификация. Това е съществено за дипломната работа, защото поставя надграждаемостта в дисциплиниран инженерен процес.

\subsection{Задължителни елементи на минимален конектор}

На практика един backend се счита за минимално интегриран, ако:

\begin{enumerate}[label=\arabic*.]
    \item съществува тип, който представлява backend connection wrapper;
    \item съществува специализация \code{connector\_trait<DB>} за този тип;
    \item специализацията дефинира \code{template <typename T> struct wire\_type};
    \item специализацията дефинира \code{cursor\_type};
    \item специализацията може да изпълнява поне базовите ORM query форми, които твърди, че поддържа.
\end{enumerate}

Това означава, че ``конектор'' не е просто рендерираща функция, а пълно описание на начина, по който C++ типовете и query IR-ът преминават към конкретния backend. Ако липсва дори един от тези елементи, добавеният backend не е архитектурно пълноценен, независимо дали част от заявките могат да бъдат изкуствено изпълнени.

\begin{table}[H]
\centering
\small
\caption{Задължителни и опционални елементи на конектора}
\label{tab:connector_requirements}
\begin{tabularx}{\textwidth}{L{4.7cm}Y}
\toprule
\textbf{Елемент} & \textbf{Роля} \\
\midrule
\code{DB} тип & wrapper над backend handle, session или mock connection \\[0.5em]
\code{connector\_trait<DB>} & централна специализация на contract-а \\[0.5em]
\code{template <typename T>\newline struct\ wire\_type} & преобразува C++ тип към wire представяне \\[0.5em]
\code{cursor\_type} & описва начина на обхождане или lifecycle-а на резултатите \\[0.5em]
\code{execute(...)} overload-и & материализират query IR-а към backend поведение \\[0.5em]
Capability tags & ограничават допустимите операции и lifecycle политики \\[0.5em]
\code{begin}/\code{commit}/\code{rollback} & необходими при transaction support \\
\code{ddl\_for(...)} & необходимо при schema-aware интеграция и \code{migrate<DB>} \\
\code{render\_constexpr(...)} & опционална compile-time materialization hook за специализирани сценарии \\
\bottomrule
\end{tabularx}
\end{table}

\subsection{Типизиране, signatures и именуване}

Connector contract-ът е строг и по отношение на типовете. \code{wire\_type<T>::type} трябва да е стабилна compile-time функция на \code{T}. \code{cursor\_type} трябва да моделира реалистичен result iteration или поне ясно да обозначава mock поведение. \code{execute(...)} overload-ите трябва да приемат \code{DB\&} като първи аргумент и да връщат \code{orm::result<...>} за \code{SELECT} форми или \code{result<std::tuple<>>} за write операции.

Именуването на конекторите в хранилището показва последователна конвенция. ``Локални'' или unit-friendly конектори използват имена като \code{SQLiteDB}, \code{PostgreSQLDB}, \code{MySQLDB}, \code{MongoDB}, \code{RedisDB}, \code{Neo4jDB} и \code{CassandraDB}. Реалните live варианти се разграничават експлицитно чрез суфикса \code{LiveDB}. Това не е дребен стилов избор, а част от contract clarity: още от името е ясно дали даден тип е mockable adapter, in-memory/local wrapper или реален driver-backed backend.

Липсата на наследяване също трябва да се разглежда като формално правило, а не като implementation detail. Нов connector не трябва да наследява ``обща ORM база'', защото това би смесило backend-специфичната логика с runtime polymorphism. Правилният механизъм за разширение е специализацията на \code{connector\_trait<DB>}.

\subsection{Capability механизъм и compile-time ограничения}

Capability tags са основният начин библиотеката да различава backend-ите по допустими операции. Ако \code{connector\_trait<DB>} не декларира \code{supports\_joins}, тогава опит за join-базирана заявка е архитектурно неправилен и следва да бъде отхвърлен още при компилация. Същото важи за транзакции, aggregation и concurrent execution.

От гледна точка на верификацията това е значимо, защото correctness вече не зависи само от тестовете. Част от contract-а е проверим директно от компилатора. Тестовете тогава не доказват съществуването на contract-а, а че конкретните специализации го прилагат последователно. Именно по тази логика \path{tests/unit/test_mongodb_connector.cpp} и \path{tests/unit/test_redis_connector.cpp} валидират отсъствието на \code{supports\_joins}, \code{supports\_transactions} и \code{supports\_aggregation} при съответните connector-и.

Тази негативна верификация е също толкова важна, колкото и позитивната. Да докажеш, че backend \emph{не} поддържа дадена операция, е толкова съществено, колкото да докажеш, че поддържа друга.

\subsection{Минимален connector skeleton и operational connector skeleton}

В хранилището \code{MockDB} и \code{MockMigrateDB} са особено ценни, защото играят ролята на connector archetypes. Първият показва operational contract-а за query execution, а вторият --- как същият contract може да бъде разширен към schema-aware migration сценарии.

\begin{lstlisting}[language=C++,caption={Извадка от \code{connector\_trait<MockDB>}},label={lst:mockdb_connector}]
template <>
struct connector_trait<MockDB>
{
    using supports_joins              = void;
    using supports_transactions       = void;
    using supports_aggregation        = void;
    using supports_upsert             = void;
    using supports_bulk_insert        = void;
    using supports_concurrent_execute = void;

    template <typename T>
    struct wire_type { using type = T; };

    struct cursor_type
    {
        [[nodiscard]] bool has_next() const noexcept { return false; }
    };

    template <typename Response, typename Joins, typename Wheres,
              typename Limits, typename Groups, typename Orders>
    static auto execute(MockDB& db,
        select_query<Response, Joins, Wheres, Limits, Groups, Orders> q)
        -> result<projected_type<Response>, Response>;
};
\end{lstlisting}

Listing~\ref{lst:mockdb_connector} е показателен, защото събира в кратка форма почти всички задължителни елементи: capability профил, \code{wire\_type}, \code{cursor\_type} и query-specific \code{execute(...)} signatures. Освен това connector-ът използва отделни overload-и за \code{select\_query}, \code{insert\_query}, \code{update\_query} и \code{delete\_query}, което прави mapping-а към backend поведението експлицитен.

\begin{lstlisting}[language=C++,caption={Извадка от schema-aware разширението \code{MockMigrateDB}},label={lst:mockmigratedb_connector}]
template <>
struct connector_trait<MockMigrateDB>
{
    using supports_joins              = void;
    using supports_transactions       = void;
    using supports_aggregation        = void;
    using supports_concurrent_execute = void;

    template <typename T>
    struct wire_type { using type = T; };

    template <typename... Args>
    static auto execute(MockMigrateDB& db, Args&&... args)
    {
        return connector_trait<MockDB>::execute(
            static_cast<MockDB&>(db), std::forward<Args>(args)...);
    }

    [[nodiscard]] static std::string ddl_for(const create_table_op& op);
    [[nodiscard]] static std::string ddl_for(const add_column_op& op);
};
\end{lstlisting}

Listing~\ref{lst:mockmigratedb_connector} показва как operational connector skeleton-ът може да бъде надграден. Тук query execution логиката се делегира към \code{MockDB}, но се добавят DDL hooks и transaction hooks. Това е добър пример за progressive connector maturity модел.

\subsection{Транзакции, нишкова безопасност и prepared queries}

Пълноценният backend contract не се изчерпва с единично изпълнение на \code{SELECT}. Ако конекторът декларира \code{supports\_transactions}, от него се очаква да поддържа \code{begin}, \code{commit} и \code{rollback}, така че \code{transaction\_guard<DB>} да бъде валиден. Ако декларира \code{supports\_concurrent\_execute}, той трябва да може да участва безопасно в \code{connection\_pool<DB, N>}.

Prepared query слоят също поставя изискване за стабилност на \code{DB\&} lifetime и за консистентно bind поведение при многократно изпълнение. Следователно зрелият конектор трябва да бъде оценяван не само по това дали рендира правилен синтаксис, а и по operational contract-а около ресурси, транзакции и повторна употреба. Това е особено важно за backend-и, които имат по-сложен драйверен lifecycle от \code{MockDB}.

\subsection{Миграции и schema-aware интеграция}

Конектор, който участва в \code{migrate<DB>}, трябва да предложи и DDL extension points. В хранилището това е демонстрирано чрез \code{MockMigrateDB}. Този пример е особено полезен, защото показва как върху вече съществуващ execution backend може да се надгради schema-aware поведение. По този начин библиотеката формализира различни нива на завършеност на конектора.

Тук е важно да се отбележи, че schema-aware поддръжката не е задължителна за всеки backend. Някои конектори могат да бъдат напълно адекватни за query execution анализ, без да поддържат автоматизирани миграции. Но ако конекторът претендира за по-пълна ORM интеграция, наличието на \code{ddl\_for(...)} hooks и migration тестове вече става задължително.

\begin{figure}[H]
\centering
\begin{tikzpicture}[node distance=1.35cm and 1.15cm]
    \node[ormaccent] (lvl1) {Ниво 1\\minimal connector};
    \node[ormblock, below=of lvl1] (lvl2) {Ниво 2\\operational connector\\query execution + result contract};
    \node[ormblock, below=of lvl2] (lvl3) {Ниво 3\\transaction/concurrency aware};
    \node[ormblock, below=of lvl3] (lvl4) {Ниво 4\\schema-aware connector\\DDL + migration hooks};
    \node[ormnote, right=1.5cm of lvl3, minimum width=4.7cm] (note)
        {\code{MockDB} покрива operational maturity, а
         \code{MockMigrateDB} показва преход към schema-aware maturity.};

    \draw[ormline] (lvl1) -- (lvl2);
    \draw[ormline] (lvl2) -- (lvl3);
    \draw[ormline] (lvl3) -- (lvl4);
\end{tikzpicture}
\caption{Стълба на зрелостта за нов connector}
\label{fig:connector_maturity_ladder}
\end{figure}

Фигура~\ref{fig:connector_maturity_ladder} предлага практичен модел за оценка. Тя е полезна и за самото бъдещо развитие на библиотеката, защото позволява новите backend-и да бъдат позиционирани не бинарно като ``има/няма connector'', а според нивото на тяхната интеграция.

\subsection{Методика за добавяне на нов backend}

Добавянето на нов backend следва ясна последователност. Първо се дефинира backend типът и неговият resource wrapper. След това се реализира \code{connector\_trait<DB>} със задължителните елементи. Следват capability декларациите и правилата за рендиране или изпълнение на query IR. Накрая се добавят unit тестове за contract-а и, при възможност, live integration тестове за реалния драйвер.

Архитектурно правилният ред е важен: не се започва от runtime workaround-и, а от формализиране на compile-time границите и на mapping логиката. Ако този ред бъде нарушен, резултатът обикновено е connector, който ``работи'' за единичен happy path, но не е съвместим с общата архитектура.

\begin{table}[H]
\centering
\small
\caption{Препоръчителна verification matrix за нов connector}
\label{tab:connector_verification_matrix}
\begin{tabularx}{\textwidth}{L{3.0cm}L{3.3cm}Y}
\toprule
\textbf{Пласт на верификация} & \textbf{Какво трябва да се докаже} & \textbf{Примерен корелат в хранилището} \\
\midrule
Concept compliance & \code{is\_connector<DB>} е удовлетворен & \path{tests/unit/test_mongodb_connector.cpp} \\
Capability honesty & декларираните и недекларираните capability tags са правилни & \path{tests/unit/test_redis_connector.cpp}, \path{tests/integration/test_sqlite.cpp} \\
Execution correctness & query IR-ът се материализира коректно & \path{tests/integration/test_mockdb.cpp} и live backend тестовете \\
Resource/lifecycle correctness & transaction и concurrency hooks се държат коректно & \path{tests/unit/test_thread_safety.cpp} \\
Schema-aware correctness & DDL hooks и migration diff-ът работят коректно & \path{tests/unit/test_schema_migration.cpp} \\
\bottomrule
\end{tabularx}
\end{table}

\subsection{Критерии за „напълно имплементиран и функциониращ“ конектор}

На базата на кода и тестовете в хранилището могат да се формулират следните критерии:

\begin{enumerate}[label=\arabic*.]
    \item конекторът удовлетворява \code{is\_connector} contract-а;
    \item capability декларациите съответстват на реалното поведение;
    \item naming и typing конвенциите са съвместими с останалите конектори;
    \item ресурсният lifecycle е коректен и проверим;
    \item поддържаните ORM операции са покрити от unit тестове;
    \item при наличен live driver има integration тестове срещу реална база данни;
    \item при заявена schema-aware поддръжка са налични DDL hooks и migration тестове.
\end{enumerate}

Тези критерии са по-строги от обикновено ``има работещ адаптер'', защото изискват консистентност между архитектурно намерение, API contract и верифицирано поведение. Именно това различава ad hoc интеграцията от конектор, който наистина принадлежи към библиотеката.

\subsection{Верификация чрез тестове}

Тестовият слой в хранилището следва структурата на архитектурата. Unit тестовете за отделните конектори валидират \code{is\_connector} compliance, capability декларациите и основното рендиране. \path{tests/unit/test_thread_safety.cpp} проверява \code{connection\_pool}, \code{thread\_local\_db} и \code{transaction\_guard}. \path{tests/unit/test_schema_migration.cpp} валидира migration diff логиката. \path{tests/integration/test_mockdb.cpp} и live integration тестовете за отделни backend-и доказват, че abstraction-ът работи и при реално изпълнение.

Следователно верификацията не е концентрирана само на едно място. Тя е разпределена така, че всяко архитектурно твърдение да има поне един пряк тестов корелат. Това превръща connector extensibility-то от декларативна цел в проверима инженерна практика.

\subsection{Оценка на наличните конектори}

\code{SQLiteDB} може да се разглежда като референтна зряла реализация за релационен backend. \code{PostgreSQLDB}, \code{MySQLDB}, \code{MongoDB}, \code{RedisDB}, \code{Neo4jDB} и \code{CassandraDB} вече демонстрират важни части от contract-а и имат различна степен на live покритие. \code{MockDB} и \code{MockMigrateDB} играят особена роля като верификационни конектори: те не са само тестови двойници, а средство за формализиране на правилата на архитектурата.

Оттук следва и по-общият извод на настоящата глава: надграждаемостта на библиотеката не е добавена впоследствие, а е заложена в самия начин, по който е дефиниран connector слоят. Именно това позволява разработката да бъде продължавана с нови backend-и, без да се нарушава формалният характер на compile-time ORM архитектурата.
